\documentclass[a4paper,12pt]{article}
\usepackage[ngerman]{babel}
\usepackage[utf8]{inputenc}
\usepackage[T1]{fontenc}
\usepackage{amsmath,amssymb}
\usepackage{tikz}
\usetikzlibrary{automata,positioning,arrows.meta}

\tikzset{
    ->,>=Stealth,
    node distance=2.5cm,
    every state/.style={thick,minimum size=1cm},
    initial text={},
}

\newcommand{\methodbox}[1]{%
\medskip
\noindent\fbox{%
    \begin{minipage}{0.94\linewidth}
    \textbf{Methodik / Definition.}\par\smallskip
    \small
    #1
    \end{minipage}%
}
\medskip\par
}

\title{ETI -- Aufgabenblatt 3 \\ Hausaufgaben}
\author{}
\date{}

\begin{document}
\raggedright
\setlength{\parindent}{0pt}

{\LARGE\bfseries ETI -- Aufgabenblatt 3 \\ Hausaufgaben\par}
\vspace{2em}

\section*{Aufgabe 3.6}

\methodbox{
Ein GNFA hat zwischen je zwei Zuständen genau eine Kante, deren Beschriftung ein regulärer Ausdruck ist.
Nicht vorhandene Kanten haben Beschriftung $\emptyset$.

\smallskip
Beim Eliminieren eines Zustands $x$ ersetzt man für alle übrigen Zustände $p,q$ die Kante
\[
    \gamma(p,q)
    \quad\text{durch}\quad
    \gamma(p,q)\cup \gamma(p,x)\gamma(x,x)^*\gamma(x,q).
\]
Danach wird $x$ entfernt. Man wiederholt das, bis nur noch Start- und Endzustand übrig sind.
Die Beschriftung der Kante vom Start- zum Endzustand ist der gesuchte reguläre Ausdruck.
}

Alle nicht eingezeichneten Kanten haben Beschriftung $\emptyset$.

\medskip
\textbf{Start-GNFA.}

\begin{center}
\begin{tikzpicture}[node distance=2.5cm and 2.6cm]
    \node[state, initial] (q0) {$q_0$};
    \node[state, right=of q0] (q1) {$q_1$};
    \node[state, right=of q1] (q2) {$q_2$};
    \node[state, below=of q2] (q3) {$q_3$};
    \node[state, accepting, above=of q2] (q4) {$q_4$};

    \path (q0) edge node[above] {$\varepsilon$} (q1)
          (q1) edge node[above] {$0$} (q2)
          (q1) edge node[above,sloped] {$\varepsilon$} (q4)
          (q2) edge[bend left=18] node[below] {$1$} (q1)
          (q2) edge node[right] {$11$} (q3)
          (q3) edge[bend left=14] node[below,sloped] {$00$} (q1);
\end{tikzpicture}
\end{center}

\medskip
\textbf{Eliminiere $q_1$.}

\begin{center}
\begin{tikzpicture}[node distance=2.5cm and 2.8cm]
    \node[state, initial] (q0) {$q_0$};
    \node[state, right=of q0] (q2) {$q_2$};
    \node[state, right=of q2] (q3) {$q_3$};
    \node[state, accepting, above=of q2] (q4) {$q_4$};

    \path (q0) edge node[above] {$0$} (q2)
          (q0) edge[bend left=25] node[above,sloped] {$\varepsilon$} (q4)
          (q2) edge[loop below] node {$10$} ()
          (q2) edge[bend left=14] node[above] {$11$} (q3)
          (q2) edge node[right] {$1$} (q4)
          (q3) edge[bend left=14] node[below] {$000$} (q2)
          (q3) edge[bend right=18] node[right] {$00$} (q4);
\end{tikzpicture}
\end{center}

\medskip
\textbf{Eliminiere $q_2$.}

\begin{center}
\begin{tikzpicture}[node distance=3cm and 4.4cm]
    \node[state, initial] (q0) {$q_0$};
    \node[state, right=of q0] (q3) {$q_3$};
    \node[state, accepting, above=of q3] (q4) {$q_4$};

    \path (q0) edge node[below] {$0(10)^*11$} (q3)
          (q0) edge[bend left=24] node[above,sloped] {$\varepsilon\cup0(10)^*1$} (q4)
          (q3) edge[loop below] node {$000(10)^*11$} ()
          (q3) edge node[right] {$00\cup000(10)^*1$} (q4);
\end{tikzpicture}
\end{center}

\medskip
\textbf{Eliminiere $q_3$.}

\begin{center}
\begin{tikzpicture}[node distance=7.6cm]
    \node[state, initial] (q0) {$q_0$};
    \node[state, accepting, right=of q0] (q4) {$q_4$};

    \path (q0) edge node[above] {$\scriptstyle
        (\varepsilon\cup0(10)^*1)\cup0(10)^*11
        (000(10)^*11)^*(00\cup000(10)^*1)$} (q4);
\end{tikzpicture}
\end{center}

\[
\boxed{
R=
\varepsilon \cup 0(10)^*1
\cup
0(10)^*11
\left(000(10)^*11\right)^*
\left(00 \cup 000(10)^*1\right)
}.
\]

\section*{Aufgabe 3.7}

\methodbox{
Pumping-Lemma zum Nachweis von Nicht-Regularität:
\begin{enumerate}
    \item Annahme: Die Sprache ist regulär. Sei $p$ die Pumping-Länge.
    \item Wähle ein Wort $s$ aus der Sprache mit $|s|\ge p$.
    \item Betrachte eine beliebige Zerlegung $s=xyz$ mit $|y|>0$ und $|xy|\le p$.
    \item Zeige für ein geeignetes $i$ meistens $i=0$ oder $i=2$, dass $xy^iz$ nicht mehr in der Sprache liegt.
\end{enumerate}
Das ist ein Widerspruch. Also ist die Sprache nicht regulär.
}

Sei
\[
A=\{w\#1^k \mid w\in\{0,1\}^*,\; w \text{ codiert die Zahl } k\}.
\]
Wir zeigen mit dem Pumping-Lemma, dass $A$ nicht regulär ist.

\medskip
\textbf{Beweis.}
\begin{enumerate}
    \item Annahme: $A$ ist regulär. Sei $p$ die Pumping-Länge.

    \item Wähle
    \[
        s=10^p\#1^{2^p}.
    \]
    Das Wort $10^p$ codiert die Zahl $2^p$, also gilt $s\in A$.
    Außerdem ist $|s|\ge p$.

    \item Sei $s=xyz$ eine Zerlegung gemäß Pumping-Lemma, also
    \[
        |y|>0
        \qquad\text{und}\qquad
        |xy|\le p.
    \]
    Dann liegt $y$ vollständig im vorderen Binärteil $10^p$, also noch vor dem Zeichen $\#$.
    Setze $k=|y|\ge 1$.

    \item Pumpe mit $i=0$. Dann bleibt der rechte Teil $\#1^{2^p}$ unverändert, aber der Binärcode links von $\#$ wird kürzer.

    Falls $y$ die führende $1$ enthält, besteht der neue Binärteil nur noch aus Nullen.
    Er codiert also nicht mehr $2^p$.
    Falls $y$ nur aus Nullen besteht, wird aus $10^p$ ein Wort der Form $10^{p-k}$.
    Dieses codiert $2^{p-k}$, und wegen $k\ge 1$ gilt $2^{p-k}\ne 2^p$.

    \item In beiden Fällen codiert der linke Teil nach dem Pumpen nicht mehr die Anzahl der Einsen rechts von $\#$.
    Also gilt
    \[
        xy^0z=xz\notin A.
    \]
    Das ist ein Widerspruch zum Pumping-Lemma.
\end{enumerate}

Folglich ist $A$ nicht regulär.

\section*{Aufgabe 3.8}

\methodbox{
Das Pumping-Lemma ist nur eine notwendige Bedingung für Regularität.
Wenn eine Sprache das Pumping-Lemma erfüllt, folgt daraus noch nicht, dass sie regulär ist.

\smallskip
Für Nicht-Regularität kann man Abschlusseigenschaften benutzen:
Sind $C$ und $R$ regulär, dann ist auch $C\cap R$ regulär.
Findet man ein reguläres $R$ mit $C\cap R=B$, wobei $B$ bekannt nicht regulär ist, dann kann $C$ nicht regulär sein.
}

Wir schreiben
\[
B=\{b^m a^{r^2}\mid m\ge 1,\; r\ge 0\}.
\]
Dann ist
\[
C = B \cup L(a^*) = B\cup \{a^n\mid n\ge 0\}.
\]

\subsection*{1. $C$ erfüllt die Bedingungen des Pumping-Lemmas}

Wir zeigen direkt die Pumping-Eigenschaft. Wähle $p=1$.
Sei $w\in C$ mit $|w|\ge 1$.

\medskip
\textbf{Fall 1: $w\in L(a^*)$.}
Dann ist $w=a^n$ mit $n\ge 1$.
Wähle
\[
    x=\varepsilon,\qquad y=a,\qquad z=a^{n-1}.
\]
Dann gilt $|y|=1>0$ und $|xy|=1\le p$.
Für alle $i\ge 0$ gilt
\[
    xy^iz = a^{n-1+i}\in L(a^*)\subseteq C.
\]

\medskip
\textbf{Fall 2: $w\in B$.}
Dann ist $w=b^m a^{r^2}$ mit $m\ge 1$.
Wähle
\[
    x=\varepsilon,\qquad y=b,\qquad z=b^{m-1}a^{r^2}.
\]
Dann gilt wieder $|y|=1>0$ und $|xy|=1\le p$.
Für $i\ge 0$ gilt
\[
    xy^iz = b^{m-1+i}a^{r^2}.
\]
Falls $m-1+i\ge 1$, liegt dieses Wort in $B$.
Der einzige andere Fall ist $m=1$ und $i=0$; dann ist
\[
    xy^0z=a^{r^2}\in L(a^*)\subseteq C.
\]
Also gilt in jedem Fall $xy^iz\in C$ für alle $i\ge 0$.

Damit erfüllt $C$ alle Bedingungen des Pumping-Lemmas.

\subsection*{2. $C$ ist nicht regulär}

Aus der Aufgabenstellung wissen wir, dass $B$ nicht regulär ist.
Betrachte die reguläre Sprache
\[
    R=L(b^+a^*)=\{b^m a^n\mid m\ge 1,\; n\ge 0\}.
\]
Dann gilt
\[
    C\cap R
    =
    (B\cup L(a^*))\cap R
    =
    B.
\]
Denn $B\subseteq R$, aber $L(a^*)\cap R=\emptyset$, da Wörter aus $R$ mindestens ein $b$ enthalten.

Angenommen, $C$ wäre regulär. Dann wäre auch $C\cap R$ regulär, weil reguläre Sprachen unter Schnitt abgeschlossen sind.
Also wäre $B$ regulär. Das widerspricht der gegebenen Aussage, dass $B$ nicht regulär ist.

Folglich ist $C$ nicht regulär.

\end{document}
